Long-tailed demonstration data is a practical obstacle for vision-language-action (VLA) policies: common manipulation skills receive many demonstrations, while rare tasks may have only a handful of trajectories. Prior work on long-tailed robotic imitation learning argues that tail failures are exposed during target approaching and proposes augmenting approach trajectories. In this paper, we study a complementary failure mode: rare-task policies often fail to bind the language-conditioned target to the demonstrated action. We introduce Task-Conditioned Action Discrimination (TCAD), a training-only objective that contrasts the likelihood of an expert action under the correct instruction against counterfactual wrong-target instructions. We combine TCAD with a moderate rare-aware behavior cloning weight, yielding Rare-Balanced Task-Conditioned Action Discrimination (RBTAD). The method adds no modules, uses no generated positive trajectories, and has no inference-time cost. A diagnostic on LIBERO-Core-LT shows that a behavior-cloned baseline assigns higher or equal action likelihood to wrong-target instructions in 64.5\% of sampled pre-contact states. In our current single-seed evaluation, RBTAD reaches 40.0\% success on LIBERO-Core-LT and improves a controlled LIBERO-Core-Full comparison from 43.0\% to 50.0\%. On a newly constructed LIBERO-Spatial-LT screening split, a relation-localized delta fusion variant improves a matched baseline from 19.0\% to 25.0\%. These results suggest that long-tailed robotic imitation benefits from directly regularizing and localizing task-conditioned action boundaries rather than only rebalancing task frequencies.
